% Stand 19. Juli 2026
\label{KapitelDatenauswertung_Diskussion}
Zielsetzung dieser Arbeit war die Untersuchung, ob die betrachteten
Feinstaubmessungen Hinweise auf eine räumliche Anisotropie 
liefern und ob sich diese bei der Modellierung mittels Kriging
berücksichtigen lässt.\\
Die Untersuchungen in dieser Arbeit liefern Indizien für das Vorliegen einer 
anisotropen Verteilung der Feinstaub-Messungen in der 
Stadt München aber keinen nach wissenschaftlichen Maßstäben
haltbaren Beweis. Insbesondere kann die Frage nach dem
Vorliegen einer  geometrischen oder zonalen Anisotropie 
nicht geklärt werden. \\
Die Ergebnisse des Kriging unter der Anisotropie-Annahme
liefern plausible Ergebnisse.\\

Die Variogramm-Karten für den Range als 
Standard-Methode zur Detektion von räum\-li\-cher Anisotropie 
liefern unplausible Ergebnisse:
der Range wird mittels einer restringierten 
numerischen Anpassung an ein theoretisches 
Kovarianz-Modell bestimmt, wobei der Range-Parameter nach oben 
durch die die maximale, paarweise Distanz der untersuchten
Sensoren beschränkt ist, die bei etwa $21-24\,k m$ liegt.\\
Da die Schätzungen der empirischen Semivarianz vielfach 
nicht die zu erwartende Sättigung in das Sill-Plateau bei
großen Werten des Lag aufweisen
(vgl.\,z.\,B.\,Abb.\,(\ref{Abb_Datenauswertung_P25MessungenApr2026DirektionaleSemivarianzBeispielFit340}))
kommt es bei der Anpassung zu numerischen Artefakten: der 
Range der Anpassung 
wird als die maximale, paarweise Distanz geschätzt.\\
%Die Untersuchung von \cite{Jang2023} oder \cite{Kim2014} deuten
% darauf hin, dass die Reichweiten mit $20-60\,km$ deutlich 
% höher liegen.\\
 
Für den Sill als einem weiteren Parameter der numerischen Anpassung
gibt es keine vorgegebene Beschränkung des Fit-Intervalls. 
Interessanterweise zeigen die Variogramm-Karten für den Sill eine bevorzugte
Richtung, in der der Sill deutlich höhere Werte annimmt als in einer 
senkrecht dazu stehenden Richtung. 
%Bei einer der drei Datensätze war es
%die Richtung $40^\circ \leftrightarrow 220^\circ$ bei den anderen beiden
%die Achse $160 \leftrightarrow 340^\circ$. 
Das ist kein stringenter 
Beweis für eine räumliche Anisotropie der Daten, könnte aber ein Indiz sein:
Sill und Range charakterisieren als $y$- und $x$-Koordinate denselben
ausgezeichneten Punkt des Übergangs vom ansteigenden
Ast hin zum Plateau, welches den Bereich der räumlichen Unkorreliertheit
der Zufallsvariablen einleitet. Dieser Zusammenhang führt zu einer 
Korrelation beider Parameter, solange 
die Beobachtungen im ansteigenden Ast der Semivarianz liegen. \\

Im letzten Teil der Arbeit werden Kriging-Vorhersagen für
anisotrop konstruierte Ko\-va\-ri\-anz-Modelle betrachtet, 
deren Parametrisierungen mit den vorherigen Beobachtungen 
konsistent sind. Diese mithilfe einer Leave-One-Out Cross-Validation
durchgeführten Analysen weisen auf eine ausgezeichnete Richtung hin,
welche in allen drei Sätzen der gemittelten Messungen gesehen wird. 
Die Aussagekraft dieser Analyse ist aufgrund der beschriebenen Einschränkungen und der hohen Varianz der
ermittelten Werte begrenzt. \\
%
%\begin{table}[H]
%\centering
%\begin{tabular}{ccc}
%\hline
%Datensatz & Vorzugsachse Sill-Karte & Vorzugsachse Kriging \\
%\hline \hline
%Juli 2025 & $160^\circ \leftrightarrow 340^\circ$ & 
%$40^\circ \leftrightarrow 220^\circ$ \\
%November 2025 & $40^\circ \leftrightarrow 220^\circ$ & 
%$40^\circ \leftrightarrow 220^\circ$ \\
%April 2026 & $60^\circ \leftrightarrow 240^\circ$ & 
%$40^\circ \leftrightarrow 220^\circ$ \\
%\end{tabular}
%\caption{Übersicht der explorativ gefundenen Vorzugsrichtungen}
%\label{Tab_Grenzwerte_Feinstaub_Landesamt}
%\end{table}

\textcolor{red}{Geometrsixhw Besonderheiten der Sensor-Positionen?}

Folglich wurden Ursachen dieser möglichen Anisotropie systematisch \textit{nicht} weiter untersucht.
Ein paar erste Vermutungen lassen sich aber anstellen:
auf Grundlage der vorliegenden Daten erscheint eine 
Erklärung anhand einer vorherrschenden Windrichtung wenig plausibel:
die betrachteten Datensätze beschreiben die monatlichen Mittelwerte von 
Sensoren, welche eine Fläche von überschlägig $20\,km \cdot 15\,km=300\,km^2$
abdecken und darüber hinaus noch unterschiedlichen Jahreszeiten zuzuordnen sind. \\
Wahrscheinlicher erscheinen zeitlich unveränderliche bauliche
oder vegetative Gegebenheiten als mögliche Erklärung.\\
% So
%würde beispielsweise der Englische Garten als eine große Grünfläche 
%mit einer ausgeprägten Längenstreckung in Richtung von etwa $58^\circ$
%eine höhere Reichweite der Korrelation des Feinstaubs entlang dieser 
%Achse erklären können (vgl.\,Abb.\,(\ref{Abb_Winkelposition_Englischer_Garten})).
%\begin{figure}[H] 
%\centering
%\includegraphics[scale = 0.44]{Abb_Winkelposition_Englischer_Garten} 
%\caption{\small{Stadtgrenzen von München und grün umrandet der Englische Garten; Zentrum der
%azimuthalen Projektion ist der Marienplatz 
%(\textcolor{gray}{$\boldsymbol{\mathrm{X}}$})}}
%\label{Abb_Winkelposition_Englischer_Garten}
%\end{figure} 


Zusammenfassend sprechen die Ergebnisse dafür, dass anisotrope Modelle bei
der räumlichen Analyse der betrachteten Feinstaubdaten eine sinnvolle
Erweiterung isotroper Modelle darstellen können. Für eine belastbare
Bestätigung wären jedoch umfangreichere Datensätze und weiterführende
 Untersuchungen erforderlich.\\

Als weitere Untersuchungen würden sich anbieten:
\begin{itemize}[topsep=2pt,parsep=1pt,itemsep=0.5pt]
\item[$\circ$] Untersuchung eines größeren Untersuchungsgebietes 
bei vergleichbarer Sensordichte.
\item[$\circ$] Die Betrachtung deutlich längerer Zeiträume,
wie sie in verwandten Arbeiten durchgeführt wurde, könnte
ebenfalls die Aussagekraft der Analysen deutlich erhöhen.
\item[$\circ$] Da die Messdaten nicht normal-verteilt sind, 
wäre es aufschlussreich, die Auswirkungen einer 
vorherigen Transformation der schiefen Daten zu testen.
\item[$\circ$] Die Untersuchung möglicher saisonaler Einflüsse 
ist ein guter Einstieg in eine mög\-licher\-wei\-se noch
detailliertere raum-zeitliche Modellierung der 
Feinstaubdaten.    
\item[$\circ$] Die Verwendung der Kriging Methode mit externem Drift 
(universal Kriging) in Kombination mit anisotropen Kovarianz-Modellen wäre 
eine weitere vielversprechenden Analyse.
\end{itemize}






